% ============================================================
\chapter{Literature Survey}
% ============================================================

This chapter surveys foundational and recent literature across the key technical
domains that PipedPeer builds upon: edge computing sandboxing and environment
portability, distributed messaging frameworks, failure-transparent distributed
systems, Python ecosystem portability, intelligent task scheduling, and
collaborative task offloading in edge networks.

% ------------------------------------------------------------------
\subsection*{Paper 1: \textit{Web Application-Based WebAssembly Container Platform
for Extreme Edge Computing}}
\addcontentsline{toc}{subsection}{Paper 1}
\textbf{Authors:} S.~Sekigawa, C.~Sasaki, and A.~Tagami \quad
\textbf{Published In:} Proc.\ IEEE GLOBECOM, 2023~\cite{sekigawa2023}

\noindent\textbf{Summary:}
Sekigawa et al.\ propose a platform that deploys WebAssembly (Wasm) modules as
isolated compute containers on extreme-edge nodes. A Wasm runtime host receives,
sandboxes, and executes user-submitted binaries without a full container daemon
such as Docker or Kubernetes. Evaluation on resource-constrained hardware shows
acceptable latency and throughput with strong isolation, demonstrating that
bytecode-level sandboxing is viable at the network edge.

\noindent\textbf{Relevance to PipedPeer:}
PipedPeer similarly executes user-submitted scripts on LAN nodes with isolated
sandboxes. The paper's insight that lightweight sandboxing can replace heavyweight
container runtimes directly motivates PipedPeer's use of \texttt{bubblewrap}
unprivileged Linux namespaces — achieving the same isolation goal without
requiring root, a container daemon, or Wasm bytecode compilation of the user's
script.

% ------------------------------------------------------------------
\subsection*{Paper 2: \textit{WebAssembly for Edge Computing: Potential and
Challenges}}
\addcontentsline{toc}{subsection}{Paper 2}
\textbf{Authors:} M.~N.~Hoque and K.~A.~Harras \quad
\textbf{Published In:} IEEE Commun.\ Standards Mag., vol.~6, no.~4, 2022~\cite{hoque2022}

\noindent\textbf{Summary:}
This survey analyses Wasm's properties — near-native speed, memory sandboxing,
and platform-agnostic bytecode — against containers and VMs for edge deployment.
Key challenges identified include WebAssembly System Interface (WASI) immaturity,
limited file-system and networking APIs, and absence of GPU support.

\noindent\textbf{Relevance to PipedPeer:}
The paper's identification of WASI's inability to support arbitrary native Python
libraries motivates PipedPeer's decision to use Nix \texttt{.nar} archives for
environment portability rather than compiling scripts to Wasm. Nix content-addressed
packages carry full native dependencies, overcoming the library-support limitations
that prevent Wasm from being a drop-in Python execution environment.

% ------------------------------------------------------------------
\subsection*{Paper 3: \textit{A Cross-Architecture Evaluation of WebAssembly in the
Cloud-Edge Continuum}}
\addcontentsline{toc}{subsection}{Paper 3}
\textbf{Authors:} S.~Kakati and M.~Brorsson \quad
\textbf{Published In:} Proc.\ IEEE/ACM CCGrid, 2024~\cite{kakati2024}

\noindent\textbf{Summary:}
This paper benchmarks Wasm runtimes (Wasmtime, WasmEdge, WAMR) across x86\_64,
ARM64, and RISC-V in cloud-edge environments. A single Wasm binary achieves
70--90\% of native performance across all architectures, with startup times under
10\,ms. The authors identify I/O-bound workloads as suffering disproportionate
overhead due to WASI syscall translation.

\noindent\textbf{Relevance to PipedPeer:}
PipedPeer must run Python workloads on workers with heterogeneous architectures
(x86\_64, aarch64, armv7l). Rather than a single platform-agnostic bytecode binary,
PipedPeer's \texttt{internal/nixgen/} generates \emph{architecture-specific} Nix
flakes for all four target platforms, ensuring full native performance and library
compatibility on each worker without the I/O overhead Wasm incurs via WASI.

% ------------------------------------------------------------------
\subsection*{Paper 4: \textit{Design of Elixir-Based Edge Server for Responsive IoT
Applications}}
\addcontentsline{toc}{subsection}{Paper 4}
\textbf{Authors:} Y.~Li and S.~Fujita \quad
\textbf{Published In:} Proc.\ IEEE CANDARW, 2023~\cite{li2023elixir}

\noindent\textbf{Summary:}
This paper builds an IoT edge server using Elixir and OTP (Open Telecom Platform).
GenServer actors handle sensor data aggregation, while supervisor trees provide
automatic process restart on failure. The system achieves sub-millisecond message
processing and maintains availability during individual process crashes.

\noindent\textbf{Relevance to PipedPeer:}
PipedPeer's daemon architecture mirrors the actor model: each node runs an
independent daemon, tasks are dispatched as messages over NATS, and
\texttt{ExecuteWithRetry} in \texttt{internal/coordinator/} implements supervisor-style
retry-on-failure. The paper's analysis of lightweight, fault-tolerant edge server
design informed PipedPeer's choice to embed all logic in a single Go daemon rather
than distributing it across microservices.

% ------------------------------------------------------------------
\subsection*{Paper 5: \textit{A Synergistic Elixir-EDA-MQTT Framework for Advanced
Smart Transportation Systems}}
\addcontentsline{toc}{subsection}{Paper 5}
\textbf{Authors:} Y.~Zhou and X.~Chen \quad
\textbf{Published In:} \textit{Computers}, vol.~16, no.~3, 2024~\cite{zhou2024}

\noindent\textbf{Summary:}
This paper combines Elixir's concurrency model, event-driven architecture (EDA), and
MQTT pub/sub messaging to build a smart transportation management system. GenStage
pipelines provide backpressure-aware event processing, and MQTT reduces communication
overhead compared to REST for high-frequency telemetry.

\noindent\textbf{Relevance to PipedPeer:}
PipedPeer uses NATS as its pub/sub transport for task dispatch and job event
publishing, analogous to MQTT in this framework. The paper's comparative analysis
of pub/sub vs.\ request-response for distributed task coordination directly supports
PipedPeer's choice of NATS over HTTP for payload delivery, and the event-driven job
lifecycle model (\texttt{internal/jobhistory/events.go}) reflects the EDA patterns
the paper evaluates.

% ------------------------------------------------------------------
\subsection*{Paper 6: \textit{WebAssembly with wasi-nn for Edge Machine Learning
Inference: Experiences and Lessons Learned}}
\addcontentsline{toc}{subsection}{Paper 6}
\textbf{Authors:} J.~Bachmeier, V.~Yussupov, J.~Henß, and H.~Koziolek \quad
\textbf{Published In:} Proc.\ ECSA, 2025~\cite{bachmeier2025}

\noindent\textbf{Summary:}
This paper evaluates Wasm with the wasi-nn API for ML inference at the edge. The
authors port image-classification and object-detection models to Wasm and benchmark
inference latency, energy consumption, and portability across runtimes. Key lessons
include wasi-nn's effective portability for CPU inference and its immaturity for GPU
acceleration.

\noindent\textbf{Relevance to PipedPeer:}
ML workloads are a primary PipedPeer use case, evidenced by the memory baselines in
\texttt{internal/resourceest/} (torch: 2\,GB, tensorflow: 3\,GB). The paper
demonstrates that Wasm is not yet viable for ML workloads requiring full CPython,
CUDA, or native libraries — validating PipedPeer's approach of shipping the complete
Python environment via Nix \texttt{.nar} archives rather than targeting a Wasm
bytecode format.

% ------------------------------------------------------------------
\subsection*{Paper 7: \textit{Case Study of WebAssembly Runtimes for AI Applications
on the Edge}}
\addcontentsline{toc}{subsection}{Paper 7}
\textbf{Authors:} S.~E.~Khelifa, M.~Bagaa, A.~O.~Messaoud, and A.~Ksentini \quad
\textbf{Published In:} 2024~\cite{khelifa2024}

\noindent\textbf{Summary:}
This case study compares Wasmtime, WasmEdge, WAMR, and Wasmer for AI inference
on edge devices. The authors measure execution time, memory overhead, and throughput
for neural network tasks, finding substantial runtime-specific variation and
identifying memory consumption as the dominant constraint on resource-limited edge
hardware.

\noindent\textbf{Relevance to PipedPeer:}
The paper's profiling of AI workload memory usage at the edge directly informs the
per-library memory baselines in PipedPeer's \texttt{internal/resourceest/} package.
The comparative runtime analysis also contextualises PipedPeer's decision to route
ML tasks to workers with sufficient free RAM using the resource estimation cascade,
rather than blindly offloading without capacity checking.

% ------------------------------------------------------------------
\subsection*{Paper 8: \textit{Programming Models for Failure-Transparent Distributed
Systems}}
\addcontentsline{toc}{subsection}{Paper 8}
\textbf{Authors:} [Fill in from DiVA portal] \quad
\textbf{Published In:} M.S.\ Thesis, KTH Royal Institute of Technology, 2025~\cite{author2025thesis}

\noindent\textbf{Summary:}
This thesis surveys programming abstractions that make distributed system failures
transparent to application developers. The author evaluates actors, reactive streams,
and process supervision trees for handling crash, omission, and Byzantine failures
through automatic retry, state replication, and task migration — without exposing
failure semantics to the application layer.

\noindent\textbf{Relevance to PipedPeer:}
PipedPeer implements a practical failure-transparent execution model:
\texttt{ExecuteWithRetry} in \texttt{internal/coordinator/} automatically migrates
tasks to the next-best worker on rejection or failure, and the hard local-fallback
guarantee ensures task completion regardless of remote node availability. The
three-phase lease protocol (accept $\to$ commit $\to$ complete) in
\texttt{internal/daemonapi/} directly reflects the thesis's framework for atomic
task reservation under partial failure.

% ------------------------------------------------------------------
\subsection*{Paper 9: \textit{Python Array API Standard: Toward Array
Interoperability in the Scientific Python Ecosystem}}
\addcontentsline{toc}{subsection}{Paper 9}
\textbf{Authors:} [Fill in from proceedings] \quad
\textbf{Published In:} Proc.\ SciPy, 2023~\cite{author2023python}

\noindent\textbf{Summary:}
This paper proposes and evaluates a standardised array computation API enabling
Python scripts to run transparently on NumPy, CuPy, JAX, and other backends without
modification. The authors demonstrate that array-interoperable scripts can be
dispatched to CPUs, GPUs, or distributed backends, supporting scientific computing
portability across heterogeneous hardware.

\noindent\textbf{Relevance to PipedPeer:}
The paper's goal — making Python scripts portable across computational backends
without modification — directly parallels PipedPeer's goal of making scripts
portable across physical nodes without code changes. PipedPeer's static import
scanner (\texttt{internal/pythondeps/}) must correctly identify array library
imports (numpy, torch, jax) and resolve them to Nix packages, ensuring that scripts
targeting the Python Array API standard are correctly packaged for remote execution.

% ------------------------------------------------------------------
\subsection*{Paper 10: \textit{iDDS: Intelligent Distributed Dispatch and Scheduling
for Workflow Orchestration}}
\addcontentsline{toc}{subsection}{Paper 10}
\textbf{Authors:} [Fill in from arXiv] \quad
\textbf{Published In:} arXiv:2510.02930, 2025~\cite{author2025idds}

\noindent\textbf{Summary:}
iDDS proposes an intelligent workflow orchestration system that dynamically assigns
scientific computing tasks to heterogeneous resources based on real-time CPU,
memory, data locality, and task-dependency constraints. A graph-based task DAG
representation drives a dynamic scheduling engine that maximises throughput while
minimising data movement.

\noindent\textbf{Relevance to PipedPeer:}
iDDS addresses the same core problem as PipedPeer — intelligent dynamic assignment
of compute tasks to heterogeneous nodes — but targets large-scale scientific
workflows rather than LAN-local script execution. The paper's scoring-based worker
selection model directly parallels PipedPeer's HealthScore dispatch model.
PipedPeer can be understood as a zero-configuration, single-binary specialisation
of this paradigm for small-team and LAN deployments.

% ------------------------------------------------------------------
\subsection*{Paper 11: \textit{Task Offloading in Multi-Hop Relay-Aided Multi-Access
Edge Computing}}
\addcontentsline{toc}{subsection}{Paper 11}
\textbf{Authors:} Y.~Deng, Z.~Chen, X.~Chen, and Y.~Fang \quad
\textbf{Published In:} IEEE Trans.\ Veh.\ Technol., vol.~72, no.~1, 2023~\cite{deng2023}

\noindent\textbf{Summary:}
This paper studies computation offloading in multi-hop wireless relay networks
reaching MEC servers. The authors formulate a joint optimisation problem for task
assignment and relay selection, minimising task completion time subject to
transmission-rate and energy constraints, and propose an iterative algorithm with
near-optimal performance.

\noindent\textbf{Relevance to PipedPeer:}
PipedPeer's task dispatch is conceptually analogous to single-hop MEC offloading:
the sender evaluates available worker nodes and routes the task to the best fit.
The paper's latency trade-off analysis between local execution and remote offloading
directly informs PipedPeer's HealthScore model, which penalises over-loaded workers
(poor remote execution time) and excludes unreachable ones (infinite latency),
and its local fallback guarantee (prefer local when offloading is infeasible).

% ------------------------------------------------------------------
\subsection*{Paper 12: \textit{Multi-Agent Deep Reinforcement Learning Based Dynamic
Task Offloading in a D2D Mobile-Edge Computing Network}}
\addcontentsline{toc}{subsection}{Paper 12}
\textbf{Authors:} H.~He, X.~Yang, X.~Mi, H.~Shen, and X.~Liao \quad
\textbf{Published In:} \textit{Sensors}, vol.~24, no.~16, 2024~\cite{he2024}

\noindent\textbf{Summary:}
This paper proposes a multi-agent deep reinforcement learning (MADRL) framework for
dynamic computation offloading in D2D MEC networks. Each device independently learns
an offloading policy from local observations (battery, channel quality, queue length)
without a central coordinator, achieving better energy-delay trade-offs than
centralised heuristics.

\noindent\textbf{Relevance to PipedPeer:}
PipedPeer uses a deterministic HealthScore model for worker selection rather than
learned policies. This paper provides a theoretical reference for what an optimal
adaptive policy could achieve, contextualising the design trade-off: PipedPeer
prioritises zero-configuration and zero training-data overhead over optimal
scheduling, accepting a simpler rule-based model that is immediately deployable
without historical data collection.

% ------------------------------------------------------------------
\subsection*{Paper 13: \textit{Joint Optimization of Multi-User Partial Offloading
Strategy and Resource Allocation in D2D-Enabled MEC}}
\addcontentsline{toc}{subsection}{Paper 13}
\textbf{Authors:} D.~Yong, R.~Liu, X.~Jia, and Y.~Gu \quad
\textbf{Published In:} \textit{Sensors}, vol.~23, no.~5, 2023~\cite{yong2023}

\noindent\textbf{Summary:}
This paper jointly optimises task partition ratios (local vs.\ offloaded), transmission
power, and resource allocation in a D2D-assisted MEC network using Lyapunov
optimisation. Partial offloading consistently outperforms all-or-nothing strategies,
with energy savings of up to 35\%.

\noindent\textbf{Relevance to PipedPeer:}
PipedPeer currently uses all-or-nothing dispatch — a script runs entirely on one
worker. The paper's partial offloading analysis motivates a direct future extension:
splitting a script's data processing stages across multiple workers and merging
results at the sender. The resource allocation model also informed PipedPeer's
multi-tier memory estimation, which matches tasks to workers with sufficient RAM
rather than blindly offloading.

% ------------------------------------------------------------------
\subsection*{Paper 14: \textit{Energy-Efficient Collaborative Task Offloading in
Multi-Access Edge Computing Based on Deep Reinforcement Learning}}
\addcontentsline{toc}{subsection}{Paper 14}
\textbf{Authors:} [Fill in from ScienceDirect] \quad
\textbf{Published In:} \textit{Ad Hoc Networks}, vol.~169, 2025~\cite{anon2025energy}

\noindent\textbf{Summary:}
This paper proposes a collaborative task offloading scheme for MEC that minimises
system-wide energy consumption using deep reinforcement learning. Devices cooperate
by relaying tasks through each other to reach edge servers, and the DRL agent learns
cooperative policies that reduce total energy by up to 40\% compared to non-cooperative
baselines.

\noindent\textbf{Relevance to PipedPeer:}
PipedPeer's collaborative compute model — where any node acts as both sender and
worker — mirrors the cooperative offloading architecture described here. The paper's
result that cooperative offloading consistently outperforms isolated local execution
validates PipedPeer's design: evaluate all available workers before falling back to
local execution, rather than defaulting to local computation immediately.

% ------------------------------------------------------------------
\subsection*{Paper 15: \textit{Blockchain Enabled Task Offloading Based on Edge
Cooperation in the Digital Twin Vehicular Edge Network}}
\addcontentsline{toc}{subsection}{Paper 15}
\textbf{Authors:} C.~Li, Q.~Chen, M.~Chen, Z.~Su, Y.~Ding, D.~Lan, and
A.~Taherkordi \quad
\textbf{Published In:} \textit{J.\ Cloud Comput.}, vol.~12, no.~120, 2023~\cite{li2023blockchain}

\noindent\textbf{Summary:}
This paper introduces a blockchain-based trust framework for cooperative task
offloading in vehicular edge networks. Smart contracts enforce service-level
agreements between task submitters and edge workers. Digital twin technology models
network state for predictive scheduling, and the blockchain prevents free-riding by
rewarding workers for verified computation.

\noindent\textbf{Relevance to PipedPeer:}
PipedPeer currently assumes cooperative, trusted participation. The paper's
trust and incentive framework directly maps onto PipedPeer's planned token-based
incentive system (Future Work, \S\ref{sec:future}), where smart contracts could
enforce execution commitments — mirroring PipedPeer's three-phase lease protocol
(accept/commit/complete) in a trustless, economically incentivised setting across
untrusted participants.

% ------------------------------------------------------------------
\subsection*{Paper 16: \textit{Congestion-Aware Distributed Task Offloading in
Wireless Multi-Hop Networks Using Graph Neural Networks}}
\addcontentsline{toc}{subsection}{Paper 16}
\textbf{Authors:} Z.~Zhao, J.~Perazzone, G.~Verma, and S.~Segarra \quad
\textbf{Published In:} arXiv:2312.02471, 2023~\cite{zhao2023}

\noindent\textbf{Summary:}
This paper models multi-hop wireless task offloading as a graph problem and trains
a graph neural network (GNN) to produce distributed offloading policies from local
neighbourhood observations. Each node makes offloading decisions using
GNN-aggregated topology information, significantly outperforming greedy baselines in
congested networks.

\noindent\textbf{Relevance to PipedPeer:}
PipedPeer's peer discovery subsystem constructs a dynamic graph of available workers
(via mDNS, registry, and peer cache), and the coordinator currently scores workers
independently using scalar HealthScore metrics. The paper's GNN-based approach
demonstrates that graph-aware placement — accounting for network topology and
inter-node latency — can substantially improve task routing. This motivates a future
PipedPeer extension in which the coordinator incorporates graph-structured worker
representations to replace independent scalar scoring.
